% ==========================================================================
% COURS 05 -- EXERCICES
% ==========================================================================

\section{Exercices du chapitre}\label{sec:sa-exercices}

\begin{TLExercise}{Exercice 1 -- Étuve thermique}
Une étuve est commandée par un gradateur. La température est mesurée par une sonde Pt100 et comparée à une consigne.
\begin{enumerate}
  \item Identifier la consigne, la sortie, la commande et deux perturbations possibles.
  \item Proposer un schéma-blocs fonctionnel.
  \item Expliquer pourquoi une commande en boucle ouverte serait insuffisante.
  \item Préciser s'il s'agit principalement d'un asservissement ou d'une régulation.
\end{enumerate}
\end{TLExercise}

\begin{TLExercise}{Exercice 2 -- Fauteuil motorisé}
Le dossier d'un fauteuil est déplacé par un motoréducteur. Un potentiomètre mesure l'angle réel.
\begin{enumerate}
  \item Compléter la chaîne : comparateur, correcteur, motoréducteur, mécanisme, capteur.
  \item Donner les unités des signaux au niveau du comparateur.
  \item Citer un défaut de capteur susceptible de créer une erreur permanente.
  \item Expliquer l'effet d'une augmentation du gain du correcteur sur précision, rapidité et stabilité.
\end{enumerate}
\end{TLExercise}

\begin{TLExercise}{Exercice 3 -- Fonction de transfert d'un circuit RC}
Un circuit vérifie :
\[
RC\frac{du_s(t)}{dt}+u_s(t)=u_e(t).
\]
\begin{enumerate}
  \item Déterminer $H(p)=U_s(p)/U_e(p)$ pour des conditions initiales nulles.
  \item Identifier le gain statique et la constante de temps.
  \item Donner la réponse à un échelon d'amplitude $E_0$.
  \item Calculer le temps de réponse à 5\%.
\end{enumerate}
\end{TLExercise}

\begin{TLExercise}{Exercice 4 -- Moteur à courant continu}
On considère :
\[
u=Ri+L\dot{i}+K_e\omega,
\qquad
J\dot{\omega}+f\omega=K_t i,
\]
et un couple résistant nul.
\begin{enumerate}
  \item Transformer les équations dans le domaine de Laplace.
  \item Éliminer $I(p)$ et déterminer $\Omega(p)/U(p)$.
  \item Donner l'ordre du modèle et identifier ses paramètres physiques.
  \item Proposer une approximation du premier ordre lorsque l'inductance est négligeable.
\end{enumerate}
\end{TLExercise}

\begin{TLExercise}{Exercice 5 -- Identification d'un premier ordre}
Un essai indiciel de $2$ V produit une variation finale de sortie de $18$ mm. La sortie atteint $11{,}4$ mm après $0{,}42$ s.
\begin{enumerate}
  \item Identifier le gain statique $K$.
  \item Identifier la constante de temps $\tau$.
  \item Estimer le temps de réponse à 5\%.
  \item Écrire la fonction de transfert et la réponse temporelle.
\end{enumerate}
\end{TLExercise}

\begin{TLExercise}{Exercice 6 -- Second ordre pseudo-périodique}
La réponse à un échelon unitaire tend vers $2$. Le premier maximum vaut $2{,}5$ et deux maxima successifs sont séparés de $0{,}8$ s.
\begin{enumerate}
  \item Calculer le dépassement relatif.
  \item En déduire le coefficient d'amortissement $\xi$.
  \item Déterminer la pulsation amortie puis la pulsation propre $\omega_0$.
  \item Écrire une fonction de transfert canonique compatible avec l'essai.
\end{enumerate}
\end{TLExercise}

\begin{TLExercise}{Exercice 7 -- Réduction d'un schéma-blocs}
Une chaîne directe est constituée de trois blocs $K$, $1/(1+\tau p)$ et $1/p$ en série. Le retour est unitaire et négatif.
\begin{enumerate}
  \item Déterminer la FTBO.
  \item Déterminer la FTBF.
  \item Donner l'équation caractéristique.
  \item Préciser la classe du système et son erreur statique pour une consigne échelon.
\end{enumerate}
\end{TLExercise}

\begin{TLExercise}{Exercice 8 -- Perturbation de couple}
Un moteur est commandé par un correcteur $C(p)$ et un procédé $P(p)$. Une perturbation $Q(p)$ s'ajoute à l'entrée du procédé. Le retour est unitaire.
\begin{enumerate}
  \item Déterminer $S(p)/E(p)$ en annulant $Q(p)$.
  \item Déterminer $S(p)/Q(p)$ en annulant $E(p)$.
  \item Montrer comment un grand gain de boucle améliore le rejet de la perturbation.
  \item Expliquer pourquoi ce gain ne peut pas être augmenté sans limite.
\end{enumerate}
\end{TLExercise}

\begin{TLExercise}{Exercice 9 -- Erreurs statiques}
On considère un retour unitaire de boucle ouverte :
\[
L(p)=\frac{20}{p(1+0{,}2p)}.
\]
\begin{enumerate}
  \item Déterminer la classe du système.
  \item Calculer l'erreur statique pour un échelon d'amplitude 3.
  \item Calculer l'erreur de traînage pour une rampe de pente $0{,}5$.
  \item Conclure sur la précision du système.
\end{enumerate}
\end{TLExercise}

\begin{TLExercise}{Exercice 10 -- Diagramme de Bode}
On considère :
\[
H(p)=\frac{50}{p(1+0{,}1p)(1+p)}.
\]
\begin{enumerate}
  \item Mettre la fonction sous forme canonique et repérer les pulsations de cassure.
  \item Tracer les asymptotes du gain.
  \item Tracer l'évolution approchée de la phase.
  \item Donner les pentes successives du gain en dB par décade.
\end{enumerate}
\end{TLExercise}

\begin{TLExercise}{Exercice 11 -- Marges de stabilité}
Pour une boucle ouverte, la pulsation de coupure à 0 dB vaut $8$ rad/s et la phase correspondante vaut $-145^\circ$. À la pulsation où la phase vaut $-180^\circ$, le gain vaut $-9$ dB.
\begin{enumerate}
  \item Calculer la marge de phase.
  \item Calculer la marge de gain.
  \item Commenter la robustesse qualitative de la boucle.
  \item Indiquer l'effet probable d'un retard supplémentaire.
\end{enumerate}
\end{TLExercise}

\begin{TLExercise}{Exercice 12 -- Synthèse complète}
Un axe linéaire doit suivre une consigne de position. Le cahier des charges impose : erreur statique nulle sur un échelon, temps de réponse à 5\% inférieur à $0{,}8$ s et dépassement inférieur à $10\%$.
\begin{enumerate}
  \item Proposer le protocole d'essai permettant de mesurer ces trois performances.
  \item Identifier les informations nécessaires pour construire un modèle.
  \item Établir une démarche de calcul de la FTBF à partir du schéma-blocs.
  \item Indiquer les grandeurs à lire dans les domaines temporel et fréquentiel.
  \item Expliquer les compromis à respecter lors du choix du correcteur.
\end{enumerate}
\end{TLExercise}
